\section{Cấu trúc khóa luận}
\label{sec:cau_truc_khoa_luan}

Nội dung khóa luận được tổ chức thành năm chương chính, cụ thể như sau:

\begin{itemize}
    \item \textbf{Chương 1 -- Tổng quan nghiên cứu:} Trình bày bài toán phát hiện té ngã, mục tiêu, phạm vi, phương pháp tiếp cận và các đóng góp chính của đề tài.
    
    \item \textbf{Chương 2 -- Cơ sở lý thuyết và tổng quan công trình liên quan:} Trình bày các khái niệm nền tảng về phát hiện té ngã dựa trên khung xương, kiến trúc YOLOv11n-Pose, cơ chế Transformer và các công trình nghiên cứu liên quan.
    
    \item \textbf{Chương 3 -- Kiến trúc hệ thống đề xuất:} Trình bày chi tiết kiến trúc tổng thể, phương pháp trích xuất đặc trưng PIFR và mô hình Hybrid Transformer phân loại.
    
    \item \textbf{Chương 4 -- Thực nghiệm và kết quả:} Trình bày quá trình thực nghiệm, bộ dữ liệu, phương pháp đánh giá, so sánh hiệu năng với các phương pháp state-of-the-art và phân tích kết quả.
    
    \item \textbf{Chương 5 -- Kết luận và hướng phát triển:} Tổng kết các kết quả đạt được, hạn chế của đề tài và đề xuất hướng nghiên cứu trong tương lai.
\end{itemize}
